\documentclass[
    language=english,
    doctype=poetry-collection,
    institution=none,
]{../../omnilatex}

\RequirePackage{config/document-settings}

\begin{document}
\maketitle

\frontmatter

\chapter*{Dedication}

\noindent\textit{For my mother, who taught me that words\\
are just breath shaped into meaning---\\
and for my father, who showed me\\
that silence speaks louder than most.}

\chapter*{Epigraph}

\begin{quote}
\textit{``Poetry is the journal of a sea animal living on land, wanting to fly in the air.''}\\
--- Carl Sandburg
\end{quote}

\chapter*{Acknowledgments}

The poems in this collection appeared, in earlier forms, in the following publications: \textit{The Paris Review}, \textit{Poetry Magazine}, \textit{The New Yorker}, \textit{Ploughshares}, \textit{The Kenyon Review}, and \textit{AGNI}. I am grateful to the editors of these journals for their support and encouragement.

Special thanks to: Ada Lim\'{o}n for her invaluable feedback on the manuscript; Ross Gay for reading early drafts with characteristic generosity; the Bread Loaf Writers' Conference for providing a space in which these poems could begin to take shape; and my students at the University of Montana, who reminded me daily why poetry matters.

\tableofcontents

\mainmatter

\part{The Body Electric}

\chapter*{I. Morning}

\begin{center}
\begin{minipage}{0.7\textwidth}
\section*{First Light}

The alarm goes off at five-fifteen,\\
the same time every morning,\\
and I lie there for a moment\\
listening to the house breathe---\\

the furnace clicking on,\\
the pipes groaning in the walls,\\
the distant hum of the refrigerator\\
keeping vigil in the kitchen.\\

Outside, the sky is the color\\
of old newsprint, pale and soft,\\
and the trees stand in their usual\\
posture of patience, waiting\\

for the sun to crest the ridge\\
and paint them gold, waiting\\
for the day to begin in earnest,\\
for the birds to start their argument\\

about territory and food\\
and the mysterious business of being alive.\\
I get up. I put my feet\\
on the cold floor. I start the coffee.\\

This is the ritual. This is the prayer\\
I say every morning without words:\\
I am here. I am still here.\\
Let me begin again.
\end{minipage}
\end{center}

\section*{The Shower}

Water falls from the showerhead\\
like a curtain of small promises,\\
each droplet carrying its own\\
brief biography of cloud and river\\
and the long, dark journey underground.\\

I stand beneath it and let it\\
wash away the residue of sleep,\\
the dreams I can't quite remember,\\
the worry that settled in my shoulders\\
sometime around three a.m.\\

The body is a simple machine:\\
it needs water, it needs heat,\\
it needs the rough comfort of soap.\\
The mind is more complicated,\\
but for now I let the water\\
do what water does best---\\
carry things away.

\section*{Breakfast}

Two eggs, scrambled. Toast with butter.\\
Coffee, black, from the French press.\\
The newspaper, folded to the crosswords,\\
which I work in pen because I am\\
either brave or foolish, I can never tell.\\

Outside the kitchen window,\\
the neighbor's cat sits on the fence\\
with the dignity of a judge,\\
observing the street with the intensity\\
of someone who knows something\\
you don't. I raise my coffee cup\\
in a silent salute. The cat blinks,\\
unimpressed. I go back to my crosswords.
\end{center}

\chapter*{II. Afternoon}

\begin{center}
\begin{minipage}{0.7\textwidth}
\section*{The Office}

The fluorescent lights hum their\\
eternal, tuneless hymn.\\
The keyboard clacks. The mouse clicks.\\
The printer wheezes and produces\\
another document no one will read.\\

I sit at my desk and pretend\\
to be productive, which is itself\\
a kind of work, requiring\\
a specific set of skills: the ability\\
to look busy while thinking\\
about something else entirely,\\
the talent for nodding at appropriate\\
intervals during meetings,\\
the craft of writing emails that say\\
nothing in the most words possible.\\

But sometimes, in the middle\\
of a Tuesday afternoon,\\
the light hits the window just so,\\
and the dust motes dance\\
in the sunbeam like tiny planets\\
orbiting an invisible star,\\
and I remember that the world\\
is full of small, unnecessary beauties,\\
and that this is reason enough\\
to keep going.
\end{minipage}
\end{center}

\section*{Lunch}

I eat my sandwich in the park,\\
on the same bench every day,\\
the one with the commemorative plaque\\
honoring someone named Harold\\
who ``loved this park and all\\
who walked within it.''\\

Today a pigeon lands\\
beside me, bold as brass,\\
and stares at my sandwich\\
with the single-minded intensity\\
of a creature that has\\
nothing else to worry about.\\

I tear off a piece of bread\\
and toss it. The pigeon\\
snatches it from the air\\
with the precision of a surgeon\\
and retreats to a safe distance,\\
where it eats with its back to me,\\
as though I might try\\
to take it back. I don't blame it.\\
Trust is earned, not given,\\
even in the bird world.

\section*{The Commute}

The bus is late, as usual.\\
I stand at the stop with the others---\\
the woman with the red umbrella,\\
the man with the briefcase,\\
the teenager with the headphones\\
who bobs his head to music\\
I can't hear. We are a small community\\
of strangers, united only\\
by our shared impatience\\
and our common destination.\\

When the bus finally arrives,\\
we board in silence,\\
dropping our coins in the box\\
with the somber ritual of acolytes\\
performing a ceremony\\
they don't quite believe in.\\
We take our seats. We stare out\\
the windows at the city sliding past,\\
this enormous, impossible city,\\
with its buildings and its people\\
and its ten thousand lights\\
beginning to flicker on\\
as the afternoon fades\\
into the blue hour,\\
the hour between day and night\\
when the world holds its breath\\
and anything seems possible.
\end{center}

\part{The World Before Windows}

\chapter*{III. Nature}

\begin{center}
\begin{minipage}{0.7\textwidth}
\section*{The Oak}

The oak in the backyard is three hundred years old.\\
I know this because I had it checked\\
by a man with a drill and a increment borer\\
who extracted a core the size of a pencil lead\\
and counted the rings like a priest counting beads.\\

Three hundred years. It was here\\
before the house, before the street,\\
before the city itself,\\
a silent witness to everything\\
that has happened on this ground---\\
the Lenape, the Dutch, the English,\\
the revolution, the reconstruction,\\
the great fire of 1897\\
that destroyed everything\\
except this tree, which stood\\
in the middle of the flames\\
and refused to burn.\\

I put my hand on its bark\\
and feel the rough, warm surface\\
beneath my fingers, and I think\\
about time, about the way it passes\\
for trees differently than for us---\\
slowly, slowly, in rings of wood\\
and leaf and root, each year\\
a complete world, a full circle,\\
a breath held and released.
\end{minipage}
\end{center}

\section*{Rain}

It rains for three days straight.\\
Not the dramatic, theatrical rain\\
of movies and novels,\\
but the quiet, persistent rain\\
of real life, the kind that seeps\\
into your bones and settles there,\\
the kind that makes you want\\
to stay in bed with a book\\
and a cup of tea and the sound\\
of water falling on the roof\\
like a thousand small fingers\\
drumming on a table,\\
patiently, patiently,\\
waiting for you to look up\\
and notice that the world\\
has been washed clean,\\
that everything is green now,\\
brighter than it was before,\\
as though the rain\\
has turned up the volume\\
on the colors,\\
made them sing\\
a little louder.

\section*{Winter}

The first snow falls\\
on the morning of December third,\\
silent and thick,\\
covering the world in white.\\

I stand at the window\\
and watch it come down,\\
each flake a tiny, intricate\\
architecture of ice,\\
unique and unrepeatable,\\
falling through the gray air\\
like a message from another world,\\
a world where things are simple\\
and clean and quiet,\\
where the only sound is the soft\\
thud of snow\\
on the branches of the oak,\\
on the roof of the house,\\
on the shoulders of the world,\\
which shrugs it off\\
and goes back to sleep.
\end{center}

\part{The Unseen}

\chapter*{IV. Memory}

\begin{center}
\begin{minipage}{0.7\textwidth}
\section*{My Mother's Hands}

My mother's hands were small\\
and quick, capable of anything---\\
baking bread, mending clothes,\\
holding a child's face\\
as though it were made of glass.\\

I remember them most\\
from the kitchen, dusted\\
with flour, moving through\\
the air with the confidence\\
of a conductor leading\\
an orchestra of ingredients,\\
transforming butter and sugar\\
and eggs into something\\
that tasted like love\\
if love had a flavor,\\
which it does---\\
it tastes like butter\\
and vanilla and the sound\\
of a timer dinging\\
in a warm kitchen\\
on a winter afternoon.
\end{minipage}
\end{center}

\section*{The House on Maple Street}

The house is gone now.\\
Demolished in 2018 to make way\\
for a condominium complex\\
with glass walls and a lobby\\
that smells of artificial lemon.\\

But I remember it---\\
the creak of the third step,\\
the way the front door stuck\\
in the humidity of August,\\
the crack in the kitchen window\\
that looked like a river\\
on a map of a country\\
that doesn't exist.\\

I remember the backyard,\\
with its rusted swing set\\
and its garden that grew\\
tomatoes and zucchini\\
and one optimistic rose bush\\
that bloomed every June\\
in defiance of the shade.\\

I remember the basement,\\
with its concrete floor\\
and its single bare light bulb\\
and the spiders who lived\\
in the corners with the dignity\\
of tenant farmers\\
who had been there\\
longer than anyone else\\
and saw no reason to move.

\section*{What the Photograph Doesn't Show}

The photograph shows a family\\
standing on a beach.\\
The sun is bright. The sky is blue.\\
Everyone is smiling.\\

What the photograph doesn't show\\
is the argument that happened\\
ten minutes before,\\
about directions,\\
about who was supposed to bring\\
the sunscreen, about whether\\
they should have come to this beach\\
at all instead of the other one,\\
the one with the parking lot\\
and the snack bar and the lifeguard\\
who looked like a movie star.\\

What the photograph doesn't show\\
is the way my father's smile\\
was forced, the way my mother's\\
eyes were red, the way my brother\\
and I stood as far apart as possible\\
while still technically being\\
``in the picture.''\\

What the photograph doesn't show\\
is that twenty years later,\\
this is the only photograph\\
I have of all four of us together,\\
and I keep it on my desk\\
because it's better to remember\\
the smile than the argument,\\
better to remember the beach\\
than the drive home\\
in silence.
\end{center}

\part{The Beautiful}

\chapter*{V. Love}

\begin{center}
\begin{minipage}{0.7\textwidth}
\section*{First Date}

You wore a blue dress.\\
I wore my best shirt,\\
which was not very best\\
but was the best I had.\\

We went to the Italian place\\
on the corner, the one\\
with the checkered tablecloths\\
and the candles in wine bottles\\
and the waiter who called\\
everyone ``my friend''\\
whether he meant it or not.\\

You ordered the lasagna.\\
I ordered the chicken parmigiana.\\
We split a bottle of Chianti\\
and talked for three hours\\
about everything and nothing---\\
about books and movies\\
and the places we'd been\\
and the places we wanted to go\\
and the mysterious way\\
the light hit your hair\\
when you turned your head\\
and the way I knew,\\
with a certainty that surprised me,\\
that this was the beginning\\
of something.
\end{minipage}
\end{center}

\section*{The Morning After the Fight}

We don't speak for two days.\\
The silence in the house\\
is a physical thing,\\
a wall between us\\
that neither of us knows\\
how to dismantle.\\

I sleep on the couch.\\
You sleep in the bedroom.\\
The door stays closed.\\

On the morning of the third day,\\
I find a note on the kitchen counter\\
in your handwriting:\\
``I'm sorry. Coffee's ready.''\\

The coffee is hot. The mug\\
is the one with the chip on the handle,\\
the one you bought at the flea market\\
because you liked the color.\\

I drink the coffee. I go to the bedroom.\\
You're sitting on the bed,\\
reading a book, pretending\\
you haven't been waiting.\\

I sit down next to you.\\
You put the book down.\\
We look at each other.\\

``I'm sorry too,'' I say.\\

You nod. You take my hand.\\

The silence breaks\\
like ice in spring,\\
like a wave on the shore,\\
like the first breath\\
after a long, long hold.

\section*{After Thirty Years}

You still leave the lights on.\\
I still forget to buy milk.\\
You still sing in the shower,\\
off-key, the same three songs\\
you've been singing since 1994.\\
I still burn the toast\\
on Sunday mornings.\\

We still argue about the thermostat,\\
about whose turn it is\\
to walk the dog,\\
about whether the paint\\
in the hallway is ``oyster white''\\
or ``cream.''\\

But we've learned something\\
over these thirty years,\\
something that took us\\
longer than it should have:\\
that love is not a destination\\
but a practice, a daily act\\
of showing up, of choosing\\
each other, again and again,\\
even when it's hard,\\
especially when it's hard.\\

And tonight, as we sit\\
on the porch and watch\\
the sun go down\\
over the same street\\
we've watched it go down over\\
for three decades,\\
you take my hand\\
and squeeze it,\\
and I squeeze back,\\
and neither of us says anything,\\
because we don't need to.\\

We've said it all already.\\
We say it again every day.\\
We'll say it tomorrow.
\end{center}

\chapter*{Colophon}

This collection was typeset using the OmniLaTeX document preparation system. The body text is set in Palatino, with display type in Century Schoolbook. The poem layouts were designed to balance readability with the white space that poetry requires to breathe.

First edition, June 2024.

\vspace{1cm}

\noindent\textcopyright\ 2024 by Eleanor Vance. All rights reserved.

\end{document}
